\title{Direct Preference Optimization: Your Language Model is Secretly a Reward Model}

\begin{abstract}
Although large-scale unsupervised language models (LMs) acquire extensive world knowledge and certain reasoning capabilities, exerting fine-grained control over their outputs remains challenging, primarily due to the unsupervised nature of their initial training. Current techniques to enhance the steerability of LMs typically involve collecting human annotations that indicate preferences between model outputs and then adapting the unsupervised LM to follow these preferences, most commonly through reinforcement learning from human feedback (RLHF). However, RLHF entails a multi-stage and frequently unstable pipeline: first, a reward model is trained to represent human preferences, and subsequently, the LM is further trained via reinforcement learning to optimize this learned reward function, while also constraining deviations from the base model. In this work, we propose a novel way of parameterizing the reward model within the RLHF paradigm, which allows the derivation of the optimal policy in closed form. This leads to a solution to the standard RLHF objective via a straightforward classification loss. The proposed algorithm, named \textit{Direct Preference Optimization} (DPO), is not only robust and efficient in practice but also removes the need for sampling from the language model during adaptation and reduces the reliance on extensive hyperparameter tuning. Through comprehensive experimentation, we demonstrate that DPO aligns LMs with human preferences as effectively or better than prevailing approaches. In particular, DPO surpasses PPO-based RLHF approaches in regulating the sentiment of generated outputs and achieves equivalent or higher quality responses for tasks like summarization and single-turn dialogue, while being markedly simpler to implement and train.
\end{abstract}

\section{Introduction}
Large unsupervised language models (LMs) that are trained on massive corpora exhibit remarkable and unexpected abilities~\citep{chowdhery2022palm, brown2020language, touvron2023llama,bubeck2023sparks}. Despite these advances, such models are trained on human-generated data encompassing a broad array of intentions, standards, and expertise. Not all of these are appropriate for direct emulation; for instance, while it is advantageous for an AI-based coding assistant to \textit{recognize} common coding errors to offer useful corrections, ideally, the assistant's code generation should be influenced by the subset of exceptional programming competence—however infrequently present—in the training data. Similarly, it is beneficial for a language model to possess \textit{awareness} of misconceptions held by a significant portion of the population (e.g., 50\%), but we do not want it to repeat these misconceptions as truth half the time it is asked. Thus, isolating and promoting a model’s \emph{preferred outputs and behavior} from its extensive \textit{background knowledge and capabilities} is central to developing AI that is reliable, effective, and controllable \citep{ouyang2022training}. While current approaches usually use reinforcement learning (RL) to align LMs with human values, we demonstrate that the RL-driven preference alignment goal traditionally adopted can, in fact, be realized exactly via a straightforward binary cross-entropy formulation, substantially streamlining preference learning.

Generally, present-day methodologies imbue LMs with preferred behavioral patterns through carefully curated datasets that encode human judgements about desirable qualities such as helpfulness and safety. This process of preference tuning is carried out after the LM undergoes initial unsupervised pre-training on large-scale text data. Although the most direct form of preference alignment is supervised fine-tuning using high-quality, human-labeled response examples, the predominant techniques today belong to the paradigm of reinforcement learning from human (or AI-generated) feedback (RLHF/RLAIF; \citep{christiano2017deep,bai2022constitutional}). RLHF approaches involve constructing a reward function using datasets of human preference judgements, then applying RL to train a language model policy that produces outputs with high predicted reward, while avoiding significant deviation from its pre-trained form. Despite the impressive dialogue and programming performance that RLHF confers, its training pipeline introduces considerable complexity over basic supervised methods, requiring additional LM training phases and repeated in-training sampling from the model’s policy, thereby imposing much greater computational demands.

This work presents a method for directly optimizing language models according to human preferences, bypassing the need for explicit reward modeling or reinforcement learning techniques. We introduce \textit{Direct Preference Optimization} (DPO), an algorithm that, while implicitly targeting the same objective as standard RLHF methods—reward maximization under a KL-divergence regularization—offers a straightforward implementation and training procedure. At its core, DPO updates aim to increase the log probability ratio between responses favored and disfavored by humans, but it additionally employs adaptive, per-sample importance weights. This adjustment helps mitigate the degenerative behavior that can arise from using a basic probability ratio objective. Similar to prior approaches, DPO draws on a theoretical preference framework (e.g., the Bradley-Terry model; \cite{bradley1952rankanalysis}) for quantifying the agreement between a reward function and observed preference data. Distinctively, while other techniques use this model to shape a loss for reward model training and subsequently optimize a policy over the learned rewards, DPO reformulates the preference loss directly in terms of the policy using a change of variables. Consequently, with access to a dataset containing human comparative judgments of model outputs, DPO is able to train a policy using a standard binary cross entropy loss, yielding a policy that is optimal for an implicit reward function aligned with observed preferences.

Our primary contribution lies in Direct Preference Optimization (DPO), a reinforcement learning-free method for fine-tuning language models with human preferences. Through empirical evaluation, we demonstrate that DPO matches or exceeds the effectiveness of conventional approaches, including RLHF with PPO, across various applications such as sentiment control, summarization, and conversational tasks, employing language models containing up to 6B parameters.

\section{Related Work}

Large self-supervised language models have demonstrated the ability to handle certain tasks with zero-shot \citep{radford2019language} or few-shot prompts \citep{gpt3,megatron,chowdhery2022palm}. Nevertheless, their efficacy on downstream benchmarks and the extent to which their responses align with user objectives are markedly enhanced by subsequent fine-tuning on collections of instructions paired with human-authored completions \citep{mishra-etal-2022-cross,sanh2022multitask,chung2022scaling,thoppilan2022lamda}. This process, known as `instruction-tuning', empowers large language models (LLMs) to extrapolate to unseen instructions and generally makes them more accessible and effective in practical use \citep{chung2022scaling}. 

While instruction-tuning has achieved notable progress, collecting \textit{relative} human evaluations of response quality is typically more scalable than assembling expert-curated examples. This has motivated subsequent research to further adapt LLMs using datasets reflecting human preferences, yielding improved capabilities in areas such as translation \citep{kreutzer-etal-2018-reliability}, summarization \citep{stiennon2022learning,ziegler2020finetuning}, story generation \citep{ziegler2020finetuning}, and instruction compliance \citep{ouyang2022training,ramamurthy2023is}. The common strategy involves first training a neural reward model to predict preferences from comparison data—typically formalized via preference models like the Bradley-Terry model \citep{bradley1952rankanalysis}—followed by refining the language model itself through reinforcement learning algorithms to optimize for this reward signal. Algorithms in use include REINFORCE \citep{williams1992reinforce}, proximal policy optimization (PPO; \cite{schulman2017proximal}), and related methods \citep{ramamurthy2023is}.

A related strand of research exploits LLMs tuned for instruction adherence and supplemented with human feedback to synthesize additional preference-labeled data, focusing on properties such as safety and harmlessness \citep{bai2022constitutional}. In these approaches, supervision is provided by human-authored textual guidelines rather than direct labeling. Together, these methodologies represent the intersection of efforts: one on training language models with reinforcement learning for complex objectives~\citep{Ranzato2015SequenceLT,paulus2018a,wu2018learning}, and the other on techniques for incorporating human preference data \citep{christiano2017deep,kupcsik2018learning}. 

Despite the attractiveness of using relative human preference data, the reinforcement learning-based fine-tuning of LLMs poses considerable practical difficulties; the present work introduces a theoretically sound alternative to directly optimize for relative preferences without relying on reinforcement learning.

Beyond the domain of language, the study of learning policies from preferences has been explored within both bandit frameworks and reinforcement learning paradigms, leading to a variety of methodologies. In contextual bandit scenarios, utilizing preferences or rankings over actions—rather than explicit reward signals—gives rise to what is termed the contextual dueling bandit (CDB; \cite{yue2012karmed,dudik2015contextual}). Without the presence of absolute rewards, the theoretical framework for CDBs replaces the standard notion of an optimal policy with that of a \textit{von Neumann winner}, defined as a policy whose expected probability of prevailing over any other policy meets or exceeds 50\% \citep{dudik2015contextual}. Distinctly, CDB setups generally receive preference labels in an online manner, whereas typical practice in human preference-based learning involves extracting information from a static, offline dataset consisting of action pairs annotated with preference judgments \citep{yan2022human}. Likewise, in \textit{preference-based RL} (PbRL), learning proceeds from binary preference signals provided by some unknown ‘scoring’ function instead of direct reward observations \citep{BusaFekete2014,ruiz2023dueling}. Although several PbRL algorithms have been proposed—some capable of leveraging off-policy preference datasets—the standard approach involves a two-step process: initially inferring the latent scoring (reward) function, followed by optimization with respect to it \citep{jain2013learning,BusaFekete2014,christiano2017deep,sadigh2017active,kupcsik2018learning}. In contrast, we introduce a unified policy optimization procedure that targets policy learning to directly fulfill expressed preferences.

\section{Preliminaries}\label{section:prelims}

Here we provide an overview of the RLHF procedure outlined by \citeauthor{ziegler2020finetuning} and subsequently employed in \citep{stiennon2022learning, bai2022training, ouyang2022training}. The RLHF process can be divided into three main stages: 1) supervised fine-tuning (SFT); 2) the collection of preference data and construction of a reward model; and 3) reinforcement learning-based optimization.

\textbf{SFT}: The RLHF workflow usually commences by adapting a pre-trained language model to specific target tasks (e.g., conversation, summarization) via supervised learning using high-quality data, resulting in an initial model $\pi^\text{SFT}$.

\textbf{Reward Modelling Phase}: During the subsequent stage, the SFT model generates answer pairs $(y_1, y_2)\sim \pi^\text{SFT}(y \mid x)$ in response to various prompts $x$. These pairs are then evaluated by human annotators, who select the response they prefer—denoted $y_w\succ y_l \mid x$, with $y_w$ as the more favored completion and $y_l$ as the less preferred. The selection is presumed to stem from an unobserved, underlying reward function $r^*(y, x)$. Multiple strategies exist for capturing these preference structures, with the Bradley-Terry (BT) \cite{bradley1952rankanalysis} model being one of the most prominent (and more complex ranking models such as the Plackett-Luce model \citep{plackett1975analysis, luce2012individual} are also applicable given access to richer preference data). Under the BT model, the distribution $p^*$ representing human preferences is formalized as:

Given a static dataset of pairwise comparisons $\mathcal{D}=\bigl\{x^{(i)}, y_w^{(i)}, y_l^{(i)}\bigr\}_{i=1}^N$ drawn from $p^*$, we can parameterize a reward function $r_{\phi}(x, y)$ and optimize its parameters via maximum likelihood estimation. Recasting the setup as a binary classification task yields the following negative log-likelihood loss:

\begin{equation}\label{eq:reward_model}
    \mathcal{L}_R(r_{\phi}, \mathcal{D}) = -\mathbb{E}_{(x, y_w, y_l)\sim \mathcal{D}}\bigl[\log \sigma(r_{\phi}(x, y_w)- r_{\phi}(x, y_l))\bigr]
\end{equation}

where the logistic function $\sigma$ is used. For large language models, $r_{\phi}(x, y)$ is typically initialized with weights from the SFT model $\pi^\text{SFT}(y \mid x)$, appending a linear projection layer atop the terminal transformer block to produce a scalar reward \cite{ziegler2020finetuning}. To limit reward variability, prior studies apply reward normalization to enforce $\mathbb{E}_{x,y\sim \mathcal{D}}\left[r_\phi(x, y)\right] = 0$ for every $x$.

\textbf{RL Fine-Tuning Phase}: In the reinforcement learning stage, the reward model delivers supervision for adjusting the language model. Following earlier approaches~\citep{jaques2017sequence, jaques2020human}, the training objective becomes

\begin{equation}\label{eq:RL}
\max_{\pi_{\theta}}  \mathbb{E}_{x\sim \mathcal{D}, y\sim \pi_{\theta}(y \mid x)}\bigl[r_{\phi}(x, y)\bigr] - \beta\mathbb{D}_{\textrm{KL}}\bigl[\pi_{\theta}(y\mid x)\mid \mid \pi_\text{ref}(y\mid x)\bigr],
\end{equation}

where $\beta$ modulates regularization against deviations from the reference policy $\pi_\text{ref}$, typically chosen as the original SFT model $\pi^\text{SFT}$. In actual training, the policy $\pi_\theta$ is also initialized with $\pi^\text{SFT}$. The KL-divergence regularization prevents excessive drift from the data regime where the reward function remains accurate, and helps sustain output diversity while avoiding collapse to a single high-reward output. Because text generation is non-differentiable, this objective requires reinforcement learning optimization. Standard procedures \citep{ziegler2020finetuning, stiennon2022learning, bai2022training, ouyang2022training} form a composite reward ${r(x, y) = r_{\phi}(x, y) -\beta (\log \pi_{\theta}(y\mid x) - \log \pi_\text{ref}(y\mid x))}$ and use PPO \cite{schulman2017proximal} for optimization.

\section{Direct Preference Optimization}\label{sec:DPO}

Because reinforcement learning algorithms are often computationally intensive and challenging to deploy at scale for language model tuning, our aim is to establish a direct preference-based policy optimization method. In contrast to standard RLHF pipelines, which first fit a reward function and then apply RL for policy updates, our method selects a specific reward parameterization so that the associated optimal policy can be written in closed form, sidestepping the need for iterative RL training. The central observation is that it is possible to utilize a known analytical correspondence between a reward function and its optimal policy, thereby reformulating reward-based loss functions into policy-based loss functions. This change-of-variable formulation eliminates the necessity for an explicit reward model, while still adhering to prevalent human preference models such as the Bradley-Terry framework. Consequently, the policy model jointly encodes the behavior of both the generative language model and the implicit reward structure.

\textbf{Derivation of the DPO objective.} Consider the RL objective as formulated in Eq.~\ref{eq:RL}, generalized by a reward function $r$, consistent with prior approaches. Based on the framework established by earlier studies~\citep{peters2007reinforcement, peng2019advantage, korbak2022reinforcement, go2023aligning}, the solution that optimizes the KL-regularized reward objective from Eq.~\ref{eq:RL} is known to be:

\begin{equation}\label{eq:op_policy}
    \pi_r(y\mid x) = \frac{1}{Z(x)}\pi_\text{ref}(y\mid x)\exp\left(\frac{1}{\beta}r(x, y)\right),
\end{equation}

where the normalization constant $Z(x) =\sum_{y}\pi_\text{ref}(y\mid x)\exp\left(\frac{1}{\beta}r(x, y)\right)$ acts as the partition function (see Appendix \ref{app:derivation1} for a full derivation). Although one could attempt to substitute the reward function $r^*$ with an estimated version $r_{\phi}$, the computation of $Z(x)$ remains intractable in practice~\citep{korbak2022reinforcement, go2023aligning}, posing significant challenges for using this representation directly. To circumvent this, we can manipulate Eq.~\ref{eq:op_policy} and represent the reward function via the optimal policy $\pi_r$, reference policy $\pi_\text{ref}$, and the partition function $Z(\cdot)$. Taking the natural logarithm of both sides in Eq.~\ref{eq:op_policy} and reorganizing yields:

\begin{equation}\label{eq:main_eq}
    r(x,y) =\beta \log \frac{\pi_r(y\mid x)}{\pi_\text{ref}(y\mid x)} + \beta \log Z(x).
\end{equation}

Applying this reparameterization to the ground-truth reward $r^*$ and its associated optimal policy $\pi^*$, and noting that the Bradley-Terry model evaluates the difference between the rewards for two responses, namely, ${p^*(y_1 \succ y_2 \mid x) = \sigma(r^*(x, y_1) - r^*(x, y_2))}$, we can now substitute Eq.~\ref{eq:main_eq} into the preference formulation in Eq.~\ref{eq:bradley-terry}. The partition function terms then cancel, resulting in an expression for the probability of a human preferring one completion over another using solely $\pi^*$ and $\pi_\text{ref}$. Therefore, under the Bradley-Terry preference model, the optimal RLHF policy $\pi^*$ adheres to:

\begin{equation}\label{eq:objective}
    p^*(y_1\succ y_2 \mid x)=\frac{1}{1 + \exp\left(\beta \log \frac{\pi^*(y_2\mid x)}{\pi_\text{ref}(y_2\mid x)} - \beta \log \frac{\pi^*(y_1\mid x)}{\pi_\text{ref}(y_1\mid x)}\right)}
\end{equation}

A comprehensive derivation is available in Appendix~\ref{app:derivation2}. While the above uses the Bradley-Terry model, this reasoning can be similarly extended to more flexible Plackett-Luce models~\citep{plackett1975analysis, luce2012individual}, with details in Appendix~\ref{app:plackett_luce_models}.

Having reformulated the likelihood of the preference data in terms of the optimal policy instead of the reward, it becomes possible to propose a maximum likelihood framework for a parametrized policy $\pi_\theta$. Drawing a parallel with the conventional reward-based approach (see Eq.~\ref{eq:reward_model}), the policy learning objective can be cast as:

\begin{equation}\label{eq:optimum_model}
    \mathcal{L}_\text{DPO}(\pi_{\theta}; \pi_\text{ref}) = -\mathbb{E}_{(x, y_w, y_l)\sim \mathcal{D}}\left[\log \sigma \left(\beta \log \frac{\pi_{\theta}(y_w\mid x)}{\pi_\text{ref}(y_w\mid x)} - \beta \log \frac{\pi_{\theta}(y_l\mid x)}{\pi_\text{ref}(y

Through this approach, we model an implicit reward via an alternative parameterization, such that the resulting optimal policy coincides with $\pi_\theta$. Additionally, since our method can be interpreted as learning a reparametrized Bradley-Terry model, it inherits desirable theoretical guarantees, including consistency under appropriate assumptions regarding the preference data distribution \cite{bong2022generalized}. In Section~\ref{sec:theory}, we provide a more detailed discussion of these theoretical aspects and situate DPO within the broader literature.

\textbf{What is the mechanism of the DPO update?} To better understand how DPO operates, we examine the gradient of its loss function, $\mathcal{L}_\text{DPO}$. The derivative with respect to the parameters $\theta$ is given by:

\begin{multline*}\label{eq:gradient}
    \nabla_\theta \mathcal{L}_\text{DPO}(\pi_\theta;\pi_\text{ref}) = \\ -\beta\mathbb{E}_{(x, y_w, y_l) \sim \mathcal{D}} \bigg[\underbrace{\sigma(\hat{r}_\theta(x, y_l) - \hat{r}_\theta (x, y_w))}_\text{greater weight when reward ordering is incorrect}\bigg[\underbrace{\nabla_\theta\log \pi(y_w \mid x)}_\text{promote $y_w$} - \underbrace{\nabla_\theta\log\pi(y_l \mid x)}_\text{demote $y_l$}\bigg]\bigg],
\end{multline*}

where the implicit reward $\hat{r}_\theta(x, y)$ is defined as $\beta \log \frac{\pi_\theta(y \mid x)}{\pi_\text{ref}(y \mid x)}$, based on both the current language model $\pi_\theta$ and a reference model $\pi_\text{ref}$ (see Section~\ref{sec:theory} for further details). The gradient of $\mathcal{L}_\text{DPO}$ intuitively serves to increase the probability assigned to preferred completions $y_w$ while decreasing the probability for less-preferred options $y_l$. Crucially, each training sample is weighted according to the degree to which the implicit reward model $\hat{r}_\theta$ erroneously ranks the dispreferred completion higher than the preferred one, modulated by the parameter $\beta$. This scaling reflects the extent to which the ordering by the implicit reward model violates the preference data, as well as the influence of the KL constraint. Our empirical analysis demonstrates that this weighting scheme is critical; omitting the weighting coefficient, as in a naive variant of the approach, can lead to degenerate behavior in the language model (see Appendix Table~\ref{tab:unlikelihood_generations}).

\textbf{DPO Overview.} The standard DPO workflow proceeds as follows: First, for each prompt $x$, generate completions $y_1, y_2 \sim \pi_\text{ref}(\cdot \mid x)$, then annotate these samples using human preference judgments to assemble an offline dataset of preferences $\mathcal{D} = \{x^{(i)}, y_w^{(i)}, y_l^{(i)}\}_{i=1}^N$. Next, the language model $\pi_\theta$ is trained to minimize $\mathcal{L}_\text{DPO}$, given $\pi_\text{ref}$, $\mathcal{D}$, and a specified value for $\beta$.

In practical applications, it is often preferable to leverage already existing public preference datasets, avoiding the overhead of producing new samples and soliciting additional human evaluations. Because such datasets typically originate from models denoted $\pi^\text{SFT}$, we set $\pi_\text{ref} = \pi^\text{SFT}$ whenever this initialization is possible. If $\pi^\text{SFT}$ is unavailable, $\pi_\text{ref}$ is constructed by maximizing the likelihood of the preferred completions ${(x, y_w)}$, formally: ${\pi_\text{ref} = \argmax_{\pi}\mathbb{E}_{x, y_w \sim \mathcal{D}}\left[\log \pi(y_w \mid x)\right]}$. This choice serves to minimize discrepancies between the inaccessible true reference distribution and the surrogate $\pi_\text{ref}$ used in DPO. Additional implementation details and hyperparameter choices are provided in Appendix~\ref{app:implementation}.

\section{Theoretical Analysis of DPO} In what follows, we deepen our analysis of DPO, develop theoretical insight, and highlight how DPO's strengths contrast with the limitations seen in actor-critic methods used for RLHF, including PPO~\cite{schulman2017proximal}.

\label{sec:theory}

\subsection{Your Language Model Is Secretly a Reward Model} DPO streamlines the learning process by obviating the need for both explicit reward function modeling and RL-based policy optimization, instead employing a single maximum likelihood training regime. The core optimization, Eq. \ref{eq:main_eq}, is analogous to the Bradley-Terry framework where rewards are parameterized as $r^*(x, y) = \beta \log\frac{\pi^*_\theta(y \mid x)}{\pi_\text{ref}(y \mid x)}$, and we directly optimize our parameterized model $\pi_{\theta}$ in a fashion equivalent to optimizing the reward model per Eq. \ref{eq:reward_model}, modulo a change of variables. In this part, we elaborate on the theoretical underpinning of this reparameterization, demonstrate that it imposes no meaningful restrictions on the class of learned reward models, and prove that the procedure can exactly reconstruct the optimal policy. We begin by formally defining an equivalence relationship over reward functions.

\begin{definition}
Two reward functions $r(x, y)$ and $r'(x, y)$ are defined to be equivalent if and only if ${r(x, y)-r'(x, y) = f(x)}$ holds for some function $f$.
\end{definition}

One can readily verify that this definition yields an equivalence relation, thereby dividing the space of reward functions into distinct equivalence classes. We can thus articulate the following lemmas:

\begin{lemma}\label{lemma:same_prefrence} Within the Plackett-Luce family, and particularly under the Bradley-Terry model, reward functions belonging to the same equivalence class produce identical preference distributions.
\end{lemma}

\begin{lemma}\label{lemma:same_policy}
    Any two reward functions from a common equivalence class will generate the same optimal policy in the context of the constrained RL formulation.
\end{lemma}

The demonstrations are direct and are presented in Appendix \ref{app:lemma1}. The first lemma reflects a widely recognized ambiguity associated with the Plackett-Luce model family \cite{plackett1975analysis}. Due to such ambiguity, it is standard practice to enforce supplementary identifiability conditions to ensure meaningful maximum likelihood estimates, as derived from Eq. \ref{eq:reward_model} \cite{bong2022generalized}. The second lemma asserts that every reward function within a given equivalence class leads to an identical optimal policy. Consequently, for our ultimate goal, it suffices to recover any representative reward function from the optimal equivalence class. In Appendix~\ref{app:thm1}, we establish the following theorem:

\begin{theorem}\label{thm:main}
    Provided mild conditions, all equivalence classes of reward functions compatible with the Plackett-Luce (as well as the Bradley-Terry) models can be realized through the reparameterization ${r(x, y) = \beta \log \frac{\pi(y\mid x)}{\pi_\text{ref}(y\mid x)}}$ for some model $\pi(y\mid x)$ and a specified reference model $\pi_\text{ref}(y \mid x)$.
\end{theorem}

\begin{sproof}
    Let $r(x, y)$ denote an arbitrary reward function, which determines the associated optimal model $\pi_r(y \mid x)$ as per Eq. \ref{eq:op_policy}. We aim to demonstrate that a reward function within the equivalence class of $r$ may be written in the form prescribed by the reparameterization. Define the projection $f$ as follows:
\begin{equation}
    f(r; \pi_\text{ref}, \beta)(x, y) = r(x, y) - \beta\log\sum_{y}\pi_\text{ref}(y\mid x)\exp\left(\frac{1}{\beta}r(x, y)\right)
\end{equation}
Here, $f$ effectively rescales $r(x, y)$ by subtracting the log-partition function of $\pi_r$. Notably, since the normalization component depends solely on $x$, $f(r; \pi_\text{ref}, \beta)(x, y) $ remains within the equivalence class of $r(x, y)$. Substituting $r$ by the right-hand side of Eq.~\ref{eq:main_eq} (applicable to any reward function), we observe $f(r; \pi_\text{ref}, \beta)(x, y) = \beta \log \frac{\pi_r(y\mid x)}{\pi_\text{ref}(y\mid x)}$. Therefore, the operation $f$ yields a member of $r$'s equivalence class in the targeted form, confirming that the suggested reparameterization is sufficiently general to represent all relevant reward functions.
\end{sproof}

Theorem~\ref{thm:main} can also be interpreted as precisely characterizing the specific reward function chosen by the DPO reparameterization within each equivalence class; namely, the function that fulfills the following condition:

\begin{equation}\label{eq:lag_p}
     \sum_{y}\underbrace{\pi_\text{ref}(y\mid x)\exp\left(\frac{1}{\beta}r(x, y)\right)}_{=\pi(y\mid x)\text{, by Thm.~\ref{thm:main} reparameterization}} = 1,
\end{equation}

In other words, $\pi(y\mid x)$ forms a legitimate probability distribution (all probabilities are non-negative and sum to unity). Moreover, referencing Eq.~\ref{eq:op_policy}, we observe that Eq.~\ref{eq:lag_p} serves as the partition function associated with the optimal policy resulting from the reward function $r(x, y)$. The DPO method's central idea is to enforce particular constraints on the typically underdetermined Plackett-Luce class of preference models—including the Bradley-Terry model—so as to preserve the expressive capacity of the class of possible reward models, while simultaneously ensuring that the optimal policy described by Eq. \ref{eq:op_policy} remains analytically tractable for any prompt $x$.

\subsection{Instability of Actor-Critic Algorithms} 

This framework further allows us to analyze the instabilities observed in conventional actor-critic methods commonly employed in RLHF, such as PPO. Adhering to the RLHF pipeline, we focus our attention on the reinforcement learning fine-tuning stage as detailed in Section \ref{section:prelims}. Connections can be drawn to the control as inference perspective \cite{levine2018reinforcement}, particularly in the context of the constrained RL scenario outlined in \ref{eq:RL}. Considering a parametrized policy $\pi_{\theta}(y\mid x)$, we seek to minimize $\mathbb{D}_{\text{KL}}[\pi_{\theta}(y|x) \mid \mid \pi^*(y\mid x)]$, with $\pi^*$ denoting the optimal policy from Eq. \ref{eq:optimum_model}, induced via the reward $r_{\phi}(y, x)$. With some mathematical manipulation, this formulation gives rise to the following optimization criterion:

\begin{equation}\label{eq:AC}
    \max_{\pi_{\theta}}\mathbb{E}_{\pi_{\theta}(y\mid x)}\left[\underbrace{r_{\phi}(x, y) -\beta\log\sum_{y}\pi_\text{ref}(y\mid x)\exp\left(\frac{1}{\beta}r_{\phi}(x, y)\right)}_{f(r_{\phi}, \pi_\text{ref}, \beta)} - \underbrace{\beta\log\frac{\pi_{\theta}(y\mid x)}{\pi_\text{ref}(y\mid x)}}_{\text{KL}}\right]
\end{equation}

The objective considered here aligns with that of previous studies 
\citep{ziegler2020finetuning, stiennon2022learning, bai2022training, ouyang2022training}, in which the DPO-equivalent reward within the $r_{\phi}$ reward class is utilized. In this context, the normalization component within $f(r_{\phi}, \pi_\text{ref}, \beta)$ can be seen as the soft value function corresponding to the reference policy $\pi_\text{ref}$. Although this term is irrelevant for determining the optimal policy, its omission can result in a high-variance policy gradient estimator, potentially destabilizing the learning process. While a learned value function could in principle be used to model this normalization, in practice such optimization is challenging. An alternative approach employed by prior literature has been to normalize rewards using a baseline derived from human completions, effectively a Monte Carlo estimate via a single sample for the normalization. Distinctively, the DPO reparameterization defines a reward function inherently independent of any such baseline.

\section{Experiments}
Here, we empirically investigate the effectiveness of DPO in learning policies from preference data. We begin with a controlled text generation task, exploring DPO’s sample efficiency in balancing reward maximization against the KL-divergence to the reference policy, and compare its behavior to that of widely-used preference optimization algorithms like PPO. We subsequently assess DPO's scalability and robustness by applying it to larger model architectures and more complex RLHF scenarios, such as tasks in summarization and dialogue. Our findings reveal that DPO, even with minimal hyperparameter adjustment, generally matches or exceeds the performance of strong baselines including PPO-based RLHF, as well as strategies involving selection of the highest-scoring trajectory among $N$ samples according to a learned reward function. We first provide an overview of our experimental protocol; more comprehensive methodological details appear in Appendix~\ref{app:exp_details}.

\textbf{Tasks.} Our study investigates three distinct open-ended text generation tasks. In each setting, algorithms aim to learn a policy from a collection of preferences $\mathcal{D}=\bigl\{x^{(i)}, y_w^{(i)}, y_l^{(i)}\bigr\}_{i=1}^N$. In the case of \textbf{controlled sentiment generation}, $x$ corresponds to the initial segment of a movie review sampled from the IMDb dataset \cite{maas-EtAl:2011:ACL-HLT2011}; the objective for the policy is to produce a continuation $y$ that reflects positive sentiment. For systematic assessment, preference pairs of generations are \textit{constructed} with the aid of a pre-trained sentiment classifier, such that $p(\text{positive}\mid x,y_w)>p(\text{positive}\mid x,y_l)$. For the SFT approach, GPT-2-large is fine-tuned until convergence using the training partition of IMDB reviews (additional details can be found in App~\ref{app:sentiment_details}). In the \textbf{summarization} task, $x$ consists of a Reddit forum post, and the policy must produce a concise summary $y$ covering the main discussion points. Consistent with previous work, we employ the Reddit TL;DR summarization dataset \citep{volske-etal-2017-tl} as well as human preference data collected by \citeauthor{stiennon2022learning}. For SFT, we utilize a model fine-tuned on human-generated summaries of forum posts,\footnote{\url{https://huggingface.co/CarperAI/openai_summarize_tldr_sft}} using the TRLX framework \citep{leandro_von_werra_2023_7790115} for RLHF training. The human preference annotations were sourced by \citeauthor{stiennon2022learning} from generations of a different, but comparably fine-tuned, SFT model. Lastly, the \textbf{single-turn dialogue} task involves $x$ as a single human input, varying from technical science inquiries to personal advice-seeking. Here, the policy must output an informative and engaging response $y$; for this, we leverage the Anthropic Helpful and Harmless dialogue dataset \citep{bai2022training}, which contains 170,000 conversations between humans and an automated assistant. In each example, two candidate replies generated by a sizable (albeit unspecified) language model are presented at the end of the conversation, with one labeled as preferred by a human annotator. Since there is no available pre-trained SFT model in this domain, we instead fine-tune an off-the-shelf language model solely using the human-preferred responses to create the SFT model.

\textbf{Evaluation.} We adopt two complementary strategies for assessment in our experiments. To investigate how well each algorithm fulfills the constrained reward maximization objective, we measure their performance in the controlled sentiment generation context using the trade-off frontier between achieved reward and KL-divergence from the base policy. This analysis is feasible here due to access to the explicit reward function, operationalized by a sentiment classifier. Conversely, in realistic applications where the underlying reward function is unknown, we instead compare each algorithm's \textit{win rate} against a baseline policy. For this, we rely on GPT-4 as a stand-in for human evaluation, judging summary quality and helpfulness in the summarization and single-turn dialogue domains, respectively. In summarization, baseline performance is determined by the gold-standard reference summaries in the test set; for dialogue, the test set's preferred response serves as the reference. While prior research indicates that large language models can surpass established automated metrics as evaluators \citep{Chen2023ExploringTU}, we also perform a human evaluation study to substantiate the use of GPT-4 in this role, detailed in Sec.~\ref{sec:human-judgments}. Our findings reveal that GPT-4's assessments are highly consistent with human judgments, and human-GPT-4 agreement generally matches or exceeds agreement between human annotators.

\textbf{Methods.} Beyond DPO, our study encompasses a range of existing methodologies for aligning language model behavior with human preferences. As baseline comparisons, we examine zero-shot prompting for the \textbf{GPT-J} model \citep{gpt-j} in summarization and two-shot prompting for \textbf{Pythia-2.8B} \citep{biderman2023pythia} in dialogue. We further evaluate the \textbf{SFT} model as well as \textbf{Preferred-FT}, which is produced by supervised fine-tuning on the favored response $y_w$, obtained from either the SFT output (in controlled sentiment and summarization scenarios) or a general-purpose language model (in the case of single-turn dialogue). Another pseudo-supervised strategy assessed is the \textbf{Unlikelihood} approach~\citep{welleck2019neural}, which trains the model to increase likelihood on $y_w$ while explicitly decreasing likelihood for $y_l$, modulated by an optional coefficient $\alpha\in[0,1]$ on the unlikelihood penalty. Additionally, we investigate \textbf{PPO} \citep{schulman2017proximal}, utilizing a reward model inferred from preference data, and \textbf{PPO-GT}, which acts as an oracle in the controlled sentiment context by leveraging the true reward function. For PPO-GT in sentiment experiments, we experiment with two variants: an off-the-shelf implementation \cite{leandro_von_werra_2023_7790115} and a custom version that rescales rewards and tunes hyperparameters to optimize results (these adjustments are also employed for PPO trained with learned rewards). Lastly, we employ a \textbf{Best of $N$} strategy, which samples $N$ candidate outputs from the SFT model (or Preferred-FT for dialogue) and selects the response with the highest reward model score from the preference data. Although this approach can yield strong results by isolating the reward model from PPO optimization challenges, it is computationally costly, as it demands $N$ sampled completions for every test-time prompt.

\subsection{How well can DPO optimize the RLHF objective?}

Within standard RLHF methodologies, the reward maximization objective is regulated by a KL-divergence constraint to simultaneously encourage high reward attainment and limit the policy's deviation from a specified reference. Consequently, evaluating algorithms in this context necessitates assessing both the obtained reward and the extent of KL divergence; a marginal reward improvement is not advantageous if it accompanies substantially increased KL divergence. Figure~\ref{fig:frontier-tldr-main} illustrates the trade-off between reward and KL across multiple algorithms under the sentiment task. For each algorithm, we perform several training trials, systematically varying the hyperparameters that determine the degree of policy conservativeness (with target KL values of $\{3,6,9,12\}$ for PPO, $\beta$ values of $\{0.05,0.1,1,5\}$, and $\alpha$ values of $\{0.05,0.1,0.5,1\}$ for unlikelihood, and differing random seeds for preferred-FT). This experimental sweep consists of 22 runs overall. After every 100 training steps—up to convergence—we evaluate all policies on a held-out set of test prompts, measuring the mean reward given by the actual reward function as well as the average sequence-level KL\footnote{That is, the sum of the per-timestep KL-divergences.} with respect to the reference policy $\text{KL}\left(\pi\mid \mid \pi_\text{ref}\right)$. Our observations indicate that DPO establishes a distinctly superior frontier, consistently reaching the highest reward for comparably low KL values. This outcome stands out for several reasons. Most notably, both DPO and PPO address the same optimization goal; nonetheless, DPO substantially surpasses PPO in efficiency, as evidenced by its strictly superior reward/KL tradeoff. Furthermore, DPO maintains an improved frontier compared to PPO even in settings where PPO has access to the true reward signal (PPO-GT).

\subsection{Can DPO scale to real preference datasets?}
\label{sec:dpo-real-datasets}
We proceed to investigate how well DPO fine-tuning performs on real-world tasks such as summarization and single-turn dialogue. In summarization tasks, traditional automatic metrics like ROUGE have been shown to correlate poorly with actual human judgments~\citep{stiennon2022learning}. Previous studies have demonstrated that preference-based fine-tuning using PPO can produce summaries more aligned with human evaluations. Our assessment involves generating samples using various methods on the test partition of the TL;DR summarization dataset and calculating the mean win rate against the reference completions found in the test set. Sampling across all methods is conducted using temperature settings ranging from 0.0 to 1.0; corresponding win rates are displayed in Figure~\ref{fig:frontier-tldr-main} (right). DPO, PPO, and Preferred-FT all leverage the same GPT-J SFT model as their initialization point\footnote{\url{https://huggingface.co/CarperAI/openai_summarize_tldr_sft}}. Our findings reveal that DPO achieves an approximate win rate of 61\% at a sampling temperature of 0.0, surpassing PPO, which attains about 57\% at its optimal (0.0) temperature. Furthermore, DPO delivers a higher peak win rate relative to the strongest result achieved by the $N$ baseline. It should be emphasized that DPO's $\beta$ hyperparameter was not systematically optimized in these experiments, suggesting that the results may not fully capture DPO’s possible effectiveness. Additionally, DPO exhibits greater resilience to changes in sampling temperature compared to PPO, which sees its performance fall to that of the underlying GPT-J model as temperatures rise. The Preferred-FT approach does not yield notable improvement over the SFT baseline. For further comparison, we directly contrast DPO and PPO via human evaluations in Section~\ref{sec:human-judgments}, where completions from DPO at a temperature of 0.25 are favored 58\% of the time over those generated by PPO at temperature 0.

For single-turn dialogue evaluation, we assess the various approaches using the subset of the test partition from the Anthropic HH dataset \citep{bai2022training} containing a single exchange between human and assistant. For GPT-4-based assessments, we utilize the preferred responses from the test set as reference outputs to determine the win rates of different models. Since there is no established SFT baseline for this benchmark, our process begins with a pre-trained Pythia-2.8B. We then apply Preferred-FT to fine-tune a reference model on the selected completions, ensuring that the outputs remain within the reference distribution, followed by further optimization with DPO. Our comparisons also include the strongest result among 128 completions produced by Preferred-FT (the Best of $N$ metric saturates at $N=128$ for this dataset; see Appendix Figure~\ref{fig:best-of-n}) and a 2-shot prompt version of the original Pythia-2.8B model. Across optimal temperature settings for each method, DPO is found to match or surpass the alternatives. Additionally, we evaluate an RLHF agent, trained with PPO on the Anthropic HH dataset \footnote{\url{https://huggingface.co/reciprocate/ppo_hh_pythia-6B}} as released by a widely-recognized group \footnote{\url{https://github.com/CarperAI/trlx/tree/main/examples/hh}}, but are unable to achieve performance improvements over the baseline Pythia-2.8B model regardless of prompt or temperature selection. Considering our findings on TL;DR and recognizing that both strategies target the same reward signal, we treat Best of 128 as a rough surrogate for PPO performance. In summary, DPO stands out as the sole computationally tractable approach to exceed the quality of preferred completions within the Anthropic HH dataset, offering comparable or superior outcomes relative to the compute-intensive Best of 128 method. Lastly, Figure~\ref{fig:dialogue-main} demonstrates that DPO rapidly attains its peak effectiveness.

\subsection{Generalization to a new input distribution}

To further explore how PPO and DPO policies respond to distributional changes, we assess their performance—trained on the Reddit TL;DR summarization task—on a new data domain: news articles from the test set of the CNN/DailyMail dataset \citep{nallapati-etal-2016-abstractive}. Both policies were tested using their optimal sampling temperatures as established for TL;DR (0 and 0.25). Table~\ref{tab:ood} reports the results of this evaluation. To measure summarization quality, we calculated the GPT-4 win rate compared to the dataset reference summaries, using the same GPT-4 (C) prompt as in the Reddit TL;DR setup, modifying the wording from ``forum post'' to ``news article''. On this out-of-distribution dataset, DPO continues to achieve noticeably higher performance than PPO. These findings offer preliminary evidence that DPO-trained policies maintain robust generalization similar to, or better than, PPO-trained policies, even though DPO does not benefit from the supplementary unlabeled Reddit TL;DR prompts available to PPO.

\subsection{Validating GPT-4 judgments with human judgments}
\label{sec:human-judgments}
To test the trustworthiness of GPT-4’s scoring, we conducted a human evaluation using outputs from the TL;DR summarization experiment, incorporating two different GPT-4 prompt variants. The \textbf{GPT-4 (S)} (simple) prompt only asks which summary better encapsulates the key information in the source post, whereas the \textbf{GPT-4 (C)} (concise) prompt also explicitly seeks a preference for conciseness; this distinction was made because, under the \textbf{GPT-4 (S)} prompt, GPT-4 tends to favor longer and more redundant summaries than human raters do. Appendix~\ref{app:prompts} contains full prompt texts. For our study, we selected three methods with varying performance—DPO at temperature 0.25 (highest), PPO at temperature 1.0 (lowest), and SFT at temperature 0.25 (intermediate)—all compared against the PPO policy sampled greedily (its strongest performing temperature). Analysis of the results indicates that GPT-4, under both prompts, agrees with human assessments as frequently as humans agree among themselves, supporting GPT-4’s validity as a surrogate for human judgment. Due to the limited pool of human annotators, multiple judgments were obtained only for comparisons involving DPO and PPO-1. Overall, the \textbf{GPT-4 (C)} prompt yields win rates that more closely align with human preferences, so this prompt was used for the principal findings in Section~\ref{sec:dpo-real-datasets}. Appendix~\ref{app:human-study} provides further methodological information, such as details about the web interface for annotators and the composition of the volunteer pool.

\section{Discussion}
Learning from preferences presents a robust and scalable approach for developing language models that are both capable and aligned with user intent. In this work, we have proposed DPO, an intuitive training method for aligning language models using preferences without the need for reinforcement learning. Instead of adapting the preference-based learning task to a standard reinforcement learning framework just to apply conventional RL algorithms, DPO establishes a direct correspondence between language model policies and reward functions. This allows for optimizing language models to match human preferences using a straightforward cross-entropy objective, sidestepping both reinforcement learning and any accompanying loss of generality. Remarkably, DPO achieves comparable or improved performance relative to prominent RLHF approaches, such as those utilizing PPO, with minimal hyperparameter tuning. As a result, DPO significantly simplifies the process of training language models with human feedback.

\textbf{Limitations \& Future Work.} Our findings prompt a number of avenues for future research. One key question is how well DPO-trained policies perform on out-of-distribution data when contrasted with models trained on explicit reward functions. Preliminary experiments indicate that DPO’s generalization is on par with PPO-based models, but a more systematic exploration is warranted. It would also be valuable to assess whether leveraging self-labeling with the DPO policy can make effective use of unlabeled prompt data. Another question concerns the nature of reward over-optimization within the direct preference optimization framework, and whether the modest drop in performance depicted in Figure~\ref{fig:dialogue-main}-right might be attributed to this effect. Furthermore, although our experiments scale to models with up to 6B parameters, investigating the scalability of DPO to contemporary models that are several orders of magnitude larger remains an open and promising research direction. With respect to evaluation, we observed that win rates as determined by GPT-4 vary based on the prompt phrasing, suggesting that future work should examine optimal strategies for eliciting reliable judgments from automated evaluation systems. Lastly, DPO holds substantial potential for broader applications beyond language modeling, including its adaptation to the training of generative models across different domains.